
\chapter{METHODOLOGY}
\label{chap:methodology}

This chapter aims to show how the experiments were built, on what data, and under which rules, in enough detail to enable reproduction. Three things are common to every experiment, so we settle them first with the computational setup and the reproducibility rules (Section~\ref{sec:m_setup}), then the pseudospectral dataset that act both as training targets and as ground truth (Section~\ref{sec:m_dataset}), and the error metrics that score every model on one scale (Section~\ref{sec:m_metrics}). Section~\ref{sec:m_ladder} then defines the six trained models. Section~\ref{sec:m_probe} sets up the probe of spectral bias, the behaviour they all share, run on the same Boussinesq solution and the same network as the rest of the work, and Section~\ref{sec:m_config} collects the hyperparameters in one table.

% ===============================================================
\section{Computational Setup and Reproducibility}
\label{sec:m_setup}

The software behind these experiments was all written for this work. It runs on Python~3.11 with PyTorch~2.12.0 \parencite{paszke2019pytorch} built against CUDA~13.0, on a single NVIDIA GeForce GTX~1660~SUPER GPU. Figures were make Matplotlib/Plotly. We used no pre-trained weights, no third-party operator-learning library, and no external dataset, although we took much inspiration from existing work. The pseudospectral solver, the multilayer perceptron, PINNs, and each optimization loop are our own code, which is what let us instrument them for the spectral analysis and the custom visualization used in upcoming chapters. The full implementation is available at \url{https://github.com/samuelkutz/TCC}.

Each learned architecture stays close to its source. The Fourier Neural Operator follows \textcite{li2021} original implementation, the Physics-Informed Neural Operator follows \textcite{li2023pino} too, together with its released code \parencite{li2023pino_code}, and the Physics-Informed Neural Network follows \textcite{raissi2019}, but implemented from scratch with PyTorch. Where an implementation departs from its reference, we explicitly point that out .

Reproducibility rests on two conventions. A single global seed initializes the Python, NumPy, and PyTorch at the start of the pipeline, so a rerun should reproduce the reported results.

Now, let us set the experimental parameters for our models and methods. Starting with the physical model, we define the spatial and temporal domains of the Boussinesq System \ref{eq:boussinesq} to be:
\begin{equation}\label{eq:m_domain}
    x \in [-L, L], \quad L = 60, \qquad t \in [0, T], \quad T = 30.
\end{equation}
Every experiment uses the initial condition of Equation~\eqref{eq:boussinesq} at unit amplitude $A = 1$, together with the initial conditions $\eta(x,0) = \sech^2(x)$ and $u(x,0) = 0$. A unit crest sits at the centre of a domain $120$ units wide. That width is deliberate. It keeps the two counter-propagating pulses of the wave-equation limit~\eqref{eq:dalembert_eta}--\eqref{eq:dalembert_u} away from the imposed periodic boundary for the whole simulated horizon. Table~\ref{tab:m_fixed} lists the constants that hold everywhere. Each of them is restated below at the point where it first matters.

\begin{table}[htbp]
    \centering
    \caption{Fixed simulation, dataset, and evaluation constants shared by every experiment. Unless a section says otherwise, these values hold throughout.}
    \label{tab:m_fixed}
    \small
    \begin{tabular}{@{}p{9.4cm}l@{}}
        \hline
        \textbf{Parameter} & \textbf{Value} \\
        \hline
        \multicolumn{2}{@{}l}{\textit{Domain and initial condition}}\\
        Spatial half-length $L$ \; ($\Omega = [-L, L]$)       & $60$ \\
        Final time $T$ \; ($t \in [0, T]$)                    & $30$ \\
        Initial amplitude $A$                                 & $1$ \\
        Initial condition                                     & $\eta(x,0) = \sech^2(x),\ u(x,0) = 0$ \\
        \hline
        \multicolumn{2}{@{}l}{\textit{Dataset solve}}\\
        Spatial grid points $N_x$                             & $128$ \\
        Time levels ($127$ RK4 steps)                         & $128$ \\
        Spatial spacing $\Delta x = 2L / N_x$                 & $0.9375$ \\
        Temporal spacing $\Delta t = T / 127$                 & $\approx 0.2362$ \\
        \hline
        \multicolumn{2}{@{}l}{\textit{Parameter sampling}}\\
        Coupled coefficients                                  & $\alpha = \beta$ \\
        Training samples $M$                                  & $20$ \\
        Sampling grid                                         & $\operatorname{linspace}(0.1,\, 4.0,\, 20)$ \\
        \hline
        \multicolumn{2}{@{}l}{\textit{Evaluation and reproducibility}}\\
        Evaluation values $\alpha = \beta$                    & $\{0.1,\, 3.21,\, 4.2\}$ \\
        Evaluation resolutions $N_x$                          & $\{128,\, 256,\, 512\}$ \\
        Spectral snapshot resolution                          & $256$ \\
        Global seed                                           & $37$ \\
        \hline
    \end{tabular}
    \\[0.4em]
    \autoriapropria
\end{table}

% ===============================================================
\section{The Parametric Dataset}
\label{sec:m_dataset}

The reference solutions do double duty. They are the supervised targets that the data-driven models learn from, and they are the ground truth against which every model is scored, working as a high-fidelity baseline in which the pseudospectral Runge--Kutta solver of Section~\ref{sec:pseudospectral} supplies it well, since we take domains big enough to consider the Boussinesq system \ref{eq:boussinesq} a periodic domain.

We keep the parameter family one-dimensional on purpose. Following Section~\ref{sec:dataset}, the initial conditions are fixed and only the PDE coefficients vary, tied together as $\alpha = \beta$ so that the two-parameter Boussinesq family collapses onto a single dimesion of difficulty. For small values of the parameters, the dynamics stay close to linear. For larger values, nonlinearities bring richier behavior. We sample parameter-solutinon pairs uniformly at $M = 20$ points,
\begin{equation}\label{eq:param_grid}
    \alpha_i = \beta_i \in \operatorname{linspace}(0.1,\, 4.0,\, 20), \qquad i = 1, \ldots, 20,
\end{equation}
running from the near-linear value $0.1$ to a more nonlinear regime at $4.0$.

Each value in~\eqref{eq:param_grid} is handed to the solver on $N_x = 128$ spatial points and marched through $127$ Runge--Kutta steps over $[0, T]$, giving $128$ time levels. The grid spacings are
\begin{equation}\label{eq:m_spacings}
    \Delta x = \frac{2L}{N_x} = \frac{120}{128} = 0.9375,
    \qquad
    \Delta t = \frac{T}{127} = \frac{30}{127} \approx 0.2362.
\end{equation}
A solve returns the two space-time fields $\eta_i(x_j, t_n)$ and $u_i(x_j, t_n)$ on the $128 \times 128$ grid, and we collect them into the solution dataset $\mathcal{S}$ of Equation~\eqref{eq:dataset}. The fields are stored as four channels in and two out, meaning we gather, for each sample, the iniital conditions $\eta_i(x_j, 0)$ and $u_i(x_j, 0) $ for each spatial point $x_j$, and the two coefficientes $\alpha_i$ and $\beta_i$, all in the spacetime grid:
\begin{equation}\label{eq:m_tensors}
    X \in \mathbb{R}^{M \times 4 \times N_x \times N_t},
    \qquad
    Y \in \mathbb{R}^{M \times 2 \times N_x \times N_t}.
\end{equation}
The two output channels carry $\eta_i$ and $u_i$. 

From now on the panels report only the surface elevation $\eta$, since it carries the wave structure we care about and the velocity field behaves the same way.

Following standard practices of deep learning, we map every channel to a common scale with a single linear, channel-wise affine transform. The scale is fixed from quantities the problem already provides rather than from the reference solutions, so that a physics-only model borrows nothing from the dataset. The elevation and velocity channels, together with the initial-condition inputs that feed them, are scaled by the known initial amplitude $A$; the coefficient channels $\alpha$ and $\beta$ by the endpoints of the sampling range~\eqref{eq:param_grid}. Writing $[f_{\min}, f_{\max}]$ for the fixed envelope of channel $f$, namely $[-A, A]$ for the fields and $[\alpha_{\min}, \alpha_{\max}]$ for the coefficients,
\begin{equation}\label{eq:m_minmax}
    \tilde f = \frac{f - f_{\min}}{f_{\max} - f_{\min} + \varepsilon},
    \qquad \varepsilon = 10^{-12}.
\end{equation}
The operators train and predict in this normalized space, and every prediction is sent back to physical units by inverting~\eqref{eq:m_minmax} before any metric or residual is formed. Because the constants are set from the initial amplitude and the sampling range, both of which the physics-only regimes of Section~\ref{sec:m_ladder} already know, the transform carries no information about the reference solutions, and the no-data models remain data-free in the output scale as well.

% ===============================================================
\section{Error Metrics}
\label{sec:m_metrics}

A comparison is only fair if every model is measured against the same reference and scales. Each method is scored against a freshly recomputed pseudospectral reference for three parameter values for \ref{eq:boussinesq},
\begin{equation}\label{eq:m_eval_params}
    \alpha = \beta \in \{0.1,\; 3.21,\; 4.2\},
\end{equation}
one near-linear and seen in the dataset, one intermediatly nonlinear and not seen in the dataset, and another even more nonlinear and outside the dataset. The middle value $3.21$ is our default whenever a single representative parameter is wanted, e.g MLPs and PINNs.

To test the discretization invariance that operators are supposed to enjoy, each trained operator is queried without retraining on grids of resolution
\begin{equation}\label{eq:m_eval_res}
    N_x \in \{128,\; 256,\; 512\},
\end{equation}
the training resolution and its two and four times refinements, chosen to also allow FFT \parencite{trefethen2000} implementations. For each query we resample the input channels of Section~\ref{sec:m_dataset} onto the finer grid and recompute the reference at the matching resolution.

Let $\eta_{\text{true}}$ and $\eta_{\text{pred}}$ be the reference and predicted elevation on a shared space-time grid. We report three complementary numbers.

\paragraph{Time-resolved relative error} At each time level $t_n$ we take the spatial $L^2$ error and divide it by the norm of the reference at that instant,
\begin{equation}\label{eq:m_rel_time}
    \mathbb{E}(t_n) = \frac{\bigl\| \eta_{\text{pred}}(\cdot, t_n) - \eta_{\text{true}}(\cdot, t_n) \bigr\|_2}{\bigl\| \eta_{\text{true}}(\cdot, t_n) \bigr\|_2 + \varepsilon}, \qquad \varepsilon = 10^{-12}.
\end{equation}
Read against time, $\mathbb{E}(t_n)$ shows where along the trajectory the error piles up, and its spread over all time levels is summarized by a box plot. This is the quantity behind the parameter and resolution panels of Chapter~\ref{chap:results}.

\paragraph{Relative spectral error} To see which spatial frequencies are wrong, we take the spatial real DFT at a fixed time and compare amplitudes mode by mode. Writing $\hat\eta(k, t_n)$ for the amplitude at wavenumber index $k$, the per-mode error is
\begin{equation}\label{eq:m_rel_spec}
    \mathbb{E}_{\text{spec}}(k, t_n) = \frac{\bigl| \, |\hat\eta_{\text{pred}}(k,t_n)| - |\hat\eta_{\text{true}}(k,t_n)| \, \bigr|}{|\hat\eta_{\text{true}}(k,t_n)|}.
\end{equation}

\paragraph{Frequency-band error during training} To watch spectral bias form as the optimizer runs, we split the spatial spectrum into three bands and save each band error every $500$ optimizer iterations. The per-mode quantity is exactly the relative spectral error~\eqref{eq:m_rel_spec}, and the band error is its average over the time levels,
\begin{equation}\label{eq:m_band_err}
    \mathbb{E}_{\text{band}}(k) = \big\langle\, \mathbb{E}_{\text{spec}}(k, t_n) \,\big\rangle_{t_n},
\end{equation}
so the same metric drives the per-frequency spectra and the band-tracking curves. With $N_k = \lfloor N_x/2 \rfloor + 1$ retained real-DFT modes, the low, mid, and high bands cut the index range into quarters,
\begin{equation}\label{eq:m_bands}
    k < \tfrac{N_k}{4}, \qquad \tfrac{N_k}{4} \le k < \tfrac{N_k}{2}, \qquad k \ge \tfrac{N_k}{2},
\end{equation}
and each band error is the average of~\eqref{eq:m_band_err} over its modes. These curves feed the band-tracking comparison of Chapter~\ref{chap:results}. Every model is tracked against one common reference, the training solution at $\alpha = \beta = 3.21$: the pointwise MLP and PINN on their own solution at that value, and each operator on the dataset sample nearest it, so the bands are read at a single parameter across all six.

% ===============================================================
\section{The Six Models}
\label{sec:m_ladder}

Four architectures give rise to six trained models, arranged along two axes. The first is the architecture family. The pointwise networks, the multilayer perceptron of Section~\ref{subsec:mlp} and the physics-informed network of Section~\ref{subsec:pinn}, represent one solution as a continuous map $(x,t) \mapsto (\eta, u)$ and are retrained for each parameter. The operator networks, the Fourier neural operator of Section~\ref{subsec:fno} and its physics-informed variant of Section~\ref{subsec:pino}, learn the whole parameter family~\eqref{eq:param_grid} at once and return a new parameter's solution in a single forward pass.

The second axis is the training regime, that is, what the loss is allowed to see:
\begin{itemize}
    \item Pure data fits the reference solutions with a supervised term alone and no equation. The data-only MLP and the FNO sit here.
    \item Data and physics keeps that supervised term and adds the Boussinesq residual beside it. The PINN and the PINO with data sit here.
    \item Pure physics switches the supervised term off, leaving the equation residual and the initial condition to drive the network. The PINN and the PINO without data sit here.
\end{itemize}
The MLP and the FNO carry no residual, so each family contributes a single pure-data model; the PINN and the PINO carry a residual whose data weight can be set on or off, so each contributes one data-and-physics and one pure-physics model. Crossing the two families with the three regimes therefore fills a $2 \times 3$ grid, the six models collected in Table~\ref{tab:m_config}.

Reading across the grid isolates what each ingredient contributes. Within a family, pure data against data and physics measures what the residual adds when the data are present, and data and physics against pure physics measures how far the equation alone carries the network once the data are removed. Across families at a fixed regime, the operator against the pointwise model measures what operator learning contributes. Every model draws on the same dataset of Section~\ref{sec:m_dataset} and is scored by the same metrics of Section~\ref{sec:m_metrics}, so a difference in behaviour traces to the family and the regime rather than to the setup.

% ===============================================================
\section{Probing the Spectral-Bias Mechanism}
\label{sec:m_probe}

In some level, all models share some spectral bias derived in Section~\ref{subsec:spectral_bias}. We examine it on the simplest object this work can offer, and we keep it on the Boussinesq system: a simple multilayer perceptron from $\mathbb{R}$ to $\mathbb{R}$ that learns one elevation profile $\eta$ at fixed time $t=15$, the with solution coming from pseudospectral again at intermediate parameter $\alpha = \beta = 3.21$,
\begin{equation}\label{eq:m_probe_target}
    x \longmapsto \eta_\theta(x) \approx \eta(x,\, 15), \qquad x \in [-L, L].
\end{equation}
One input, one output, one field, one instant. Nothing is introduced that the rest of the work does not already use: no auxiliary equation, no synthetic target, no second architecture. Serving as a proof of concept of the theory that Section~\ref{subsec:spectral_bias} speaks of, a target function sampled on a regular periodic grid and fit by gradient descent.

Two choices deserve a word. Fixing the time removes the temporal axis from the problem, so the Fourier analysis of the error is a plain spatial transform on the $N_x = 128$ grid and the wavenumbers of the theory are the wavenumbers of the figure, with no averaging in between. The instant $t = 15$ sits at the middle of the horizon.

\subsection{The Probe Setting}
\label{sec:m_probe_setting}

We train three networks, at widths $N_{\mathrm{MLP}} \in \{8, 64, 512\}$, each with four hidden layers, hyperbolic-tangent activations, one input, one output, and the spatial coordinate mapped to $[-1,1]$, before the forward pass. Every figure in this section comes from those same three runs.

Because the target lives on a single grid of $N_x = 128$ points, the Jacobian $J \in \mathbb{R}^{N_x \times N_\theta}$ is cheap enough to assemble explicitly, one backward pass per grid point, at every width including the widest. That matters, because the empirical kernel
\begin{equation}\label{eq:m_probe_kernel}
    K_0 = J_0 J_0^{\mathsf{T}} = V \Lambda V^{\mathsf{T}}
\end{equation}
has to be eigendecomposed, not merely applied, and it is only $128 \times 128$. It is also cheap enough to rebuild the kernel during training, which allow us to see how much it drifts from the initial state. Every run is carried out in double precision so the small eigenvalues survive the decomposition. The three widths are trained with the same learning rate, a fixed fraction of the stability bound of Section~\ref{subsec:spectral_bias}, and for the same number of iterations, aiming fair comparisions.

Training is full-batch gradient descent on the sum-of-squares loss $\tfrac{1}{2}\|e\|^2$, so that a single step is exactly
\begin{equation}\label{eq:m_theory_step}
    \theta \leftarrow \theta - \mu\, J^{\mathsf{T}} e,
\end{equation}
the update behind~\eqref{eq:error_discrete_step}, with $e = \eta_\theta - \eta_{\text{true}}$ measured on the grid. The learning rate is fixed to the same fraction of each width's stability bound, $\mu = 0.2\,(2/\lambda_{\max}) = 0.4/\lambda_{\max}$, with $\lambda_{\max}$ the largest eigenvalue of $K_0$ at that width. Anchoring $\mu$ to the bound $\mu < 2/\lambda_{\max}$ of Section~\ref{subsec:spectral_bias} rather than to a randomly chosen value keeps the three widths comparable too. We iterate $500{,}000$ times.

One caveat has to be explained, the MLP is nonlinear in its parameters, so its Jacobian is not constant: $J$ drifts away from $J_0$ as training proceeds, and the initial-kernel prediction $e^{(n)} = (I - \mu K_0)^n e_0$ of~\eqref{eq:error_discrete} is an approximation here, coming from the gradient descent/forward euler discretization of the gradient flow \ref{eq:error_linear_ode}.

Therefore, we analyse the three predictions and compare them against the initial-kernel predictions of Section~\ref{subsec:spectral_bias}.
\begin{itemize}
    \item Error decay $e(\tau)$. We project the measured error onto each eigenvector $v_k$ of $K_0$, giving $c_k^{(n)} = v_k^{\mathsf{T}} e^{(n)}$, and compare its amplitude against the predicted $|c_k(0)|\,|1 - \mu\lambda_k|^n$ of~\eqref{eq:error_discrete}. Because the decay times $1/(\mu\lambda_k)$ span several decades, the tracked eigendirections are chosen to spread over the timescales the run can resolve, and the iteration axis is logarithmic.
    \item Spectral decay $\hat e(\tau)$. We take the real DFT of the error profile at each iteration, giving one amplitude per wavenumber, and compare the measured decay against the prediction of~\eqref{eq:fourier_evolution}. The wavenumbers are fixed at physical values, so the panel reads the frequency principle literally.
    \item Per-mode iteration count $n_k$. For each eigendirection we record the iteration at which $|c_k^{(n)}|$ first drops below a precision $\varepsilon$ and compare it against the discrete crossing time $n_k = \lceil \ln(|c_k(0)|/\varepsilon)/(-\ln|1-\mu\lambda_k|)\rceil$ of~\eqref{eq:iteration_bound_discrete}, the exact step at which the geometric law $|1-\mu\lambda_k|^n$ of~\eqref{eq:error_discrete} falls below $\varepsilon$. Using this discrete count rather than its continuous small-step limit~\eqref{eq:iteration_count} keeps the panel on the same $(1-\mu\lambda_k)^n$ law as the decay curves, so a departure from the diagonal reflects kernel drift rather than the truncation of the exponential. We set $\varepsilon = 0.05\,\max_k |c_k(0)|$, a fraction of the largest initial coefficient rather than an absolute number, so the threshold scales with the amplitude of $\eta$ instead of with the units it is measured in. The closed form carries no width, so the three networks should fall on one diagonal; both axes are logarithmic, since the counts span decades.
\end{itemize}
The first two panels are drawn at the widest network, where the initial-kernel picture has the best claim to hold; the third collects all three widths.

\subsection{Kernel Drift and the Eigenvalue Spectrum}
\label{sec:m_probe_ntk}

The predictions above assume the kernel stays put, and we measure that during the very runs whose predictions we are testing. We log the relative parameter drift $\|\theta - \theta_0\| / \|\theta_0\|$ and the relative kernel drift $\|K - K_0\| / \|K_0\|$; the eigenvalue spectrum of $K$ is recorded at initialization and again at the end of training.

The first two ask whether the kernel stays put, as lazy training assumes, and whether it does so more convincingly as the network widens, which is what the initial-kernel approximation~\eqref{eq:error_linear_ode} predicts. The third exposes the eigenvalue disparity that~\eqref{eq:training_time} names as the cause of spectral bias. Together with the three checks of Section~\ref{sec:m_probe_setting}, these curves open the results of Chapter~\ref{chap:results}, where they anchor the reading of the spectral-bias curves measured on the full models.

% ===============================================================
\section{Configuration Summary}
\label{sec:m_config}

Table~\ref{tab:m_config} gathers the configuration of every experiment. The four learned architectures share a common budget, $15{,}000$ iterations and, for the Adam-optimized models, learning rate $10^{-3}$, so the comparison of Chapter~\ref{chap:results} reflects the architecture and the training regime rather than the effort spent on tuning. The last row reports the measured training wall-clock on the hardware of Section~\ref{sec:m_setup}, which is indicative and depends on the GPU used.

\begin{table}[htbp]
    \centering
    \caption{Configuration of the four training architectures and of the spectral-bias probe. The FNO and PINO share the operator backbone and its parameter count; the MLP and PINN share the feed-forward backbone and are single-parameter models. All Adam-trained models use $15{,}000$ iterations. The probe is a separate, deliberately minimal network, $\mathbb{R} \to \mathbb{R}$ on the single profile $\eta(x, 15)$, trained once per width.}
    \label{tab:m_config}
    \footnotesize
    \setlength{\tabcolsep}{4pt}
    \begin{tabular}{lccccc}
        \hline
        \textbf{Setting} & \textbf{MLP} & \textbf{PINN} & \textbf{FNO} & \textbf{PINO} & \textbf{Bias probe} \\
        \hline
        Input $\to$ output        & $\mathbb{R}^2 \to \mathbb{R}$ & $\mathbb{R}^2 \to \mathbb{R}^2$ & operator & operator & $\mathbb{R} \to \mathbb{R}$ \\
        Trainable parameters      & $198{,}401$ & $198{,}658$ & $2{,}106{,}146$ & $2{,}106{,}146$ & $241,\ 12{,}673,\ 789{,}505$ \\
        Hidden layers             & $4$         & $4$         & $4$ Fourier     & $4$ Fourier     & $4$ \\
        Width / channels          & $256$       & $256$       & $32$            & $32$            & $8, 64, 512$ \\
        Fourier modes (per axis)  & ---         & ---         & $16$            & $16$            & --- \\
        Optimizer                 & Adam        & Adam        & Adam            & Adam            & GD \\
        Learning rate             & $10^{-3}$   & $10^{-3}$   & $10^{-3}$       & $10^{-3}$       & $0.4/\lambda_{\max}$ \\
        Iterations                & $15{,}000$  & $15{,}000$  & $15{,}000$      & $15{,}000$      & $500{,}000$ \\
        Parameters trained        & single ($3.21$) & single ($3.21$) & all $20$    & all $20$        & single ($3.21$) \\
        Collocation / batch       & full batch  & $5{,}000$ pts & batch $256$    & batch $256$    & $128$ pts (full) \\
        Loss weights pde/ic/data  & $-/-/1$     & $1/1/\{1,0\}$ & $-/-/1$        & $100/1/\{1,0\}$  & $-/-/1$ \\
        Physics-residual grid     & ---         & $5{,}000$ pts & ---            & $256 \times 256$ & --- \\
        Target                    & $\eta(x,t)$ & $\eta, u$     & family         & family           & $\eta(x, 15)$ \\
        Data regimes              & data only   & with / without & supervised   & with / without & data only \\
        Training wall-clock       & $\approx 143$\,s & $\approx 880$\,s & $\approx 1101$\,s & $\approx 4200$--$5200$\,s & --- \\
        \hline
    \end{tabular}
    \\[0.4em]
    \autoriapropria
\end{table}
